\section{Two Fireflies: An Exact Answer}
The full model may contain hundreds of phases, but the case N = 2 can be solved directly. Define
\[
\Delta = \theta_{2} - \theta_{1}, \quad \Delta \omega = \omega_{2} - \omega_{1}.
\]
For two fireflies, eq.~\ref{eqn:3} becomes
\begin{align}
    \frac{d\omega_{1}}{dt} &= \omega_{1} + \frac{K}{2} \sin (\theta_{2}-\theta_{1}) \\
    \frac{d\omega_{2}}{dt} &= \omega_{2} + \frac{K}{2} \sin (\theta_{1}-\theta_{2})
\end{align}

Subtracting the first equation from the second gives
\begin{align}
    \frac{d\Delta}{dt} &= \Delta\omega + \frac{K}{2} \sin (-\Delta) - \frac{K}{2} \sin (\Delta)\\
     &= \Delta\omega - K \sin \Delta \label{eqn:7}
\end{align}
The two fireflies are phase locked when $\Delta$ becomes constant. Setting the left side of eq.~\ref{eqn:7} to zero 
\begin{equation}
    \sin \Delta^* = \frac{\Delta\omega}{K}
\end{equation}

A real solution exists only when
\begin{equation}
    |\Delta\omega| \le K.
\end{equation}

Thus, coupling can overcome a small mismatch between internal clocks, but not an arbitrarily large one.